\scanpage{45}
\mspage{43}

Mais \unsure{on} peut toujours \uline{choisir} \deleted{$\unsure{K,K}$} un \deleted{$\unsure{\overline K}$} objet \inserted{$\overline K$} de Cat, et une flèche

\[
 i:K\longrightarrow\overline K\qquad\text{\uline{dans Cat}}
\]

pas seulement \deleted{\unsure{dans}} Cat hot, de telle façon de plus que $i$ soit une ``\uline{cofibration} de Thomason''. (et en particulier une \unsure{imm ouverte|imm. ouverte} ?). Il \deleted{s'ensuit} en résulte que si on a un système inductif

\[
 K=(K_n)_{n\in\mathbb N}
\]

indexé par $\mathbb N$, et si \deleted{\unsure{on}} on choisit pour chaque $n$ une enveloppe de Thomason-Kan

\[
 K_n\overset{\alpha_n}{\hookrightarrow}\overline K_n
\]

de cette façon, on peut trouver des \unsure{morphismes|flèches} de transition

\[
 \overline K_n\overset{\overline{\unsure{i_n|j_n}}}{\unsure{\relbar|\longrightarrow}}\overline K_{n+1}
\]

\unsure{tels} que les carrés

\diagram{\begin{tikzcd}[ampersand replacement=\&]
 K_n\arrow[r,hook,"{\alpha_n}"]\arrow[d,swap,"{i_n}"] \& \overline K_n\arrow[d,"{\overline{i_n}}"]\\
 K_{n+1}\arrow[r,hook,"{\alpha_{n+1}}"] \& \overline K_{n+1},
\end{tikzcd}}

soient commutatifs ; de telle sorte que l'on ait donc un hom de \uline{systèmes inductifs de Cat}

\[
 K\longrightarrow\overline K
\]

et a fortiori un hom de Ind-objet de Cat. Ceci posé, les conditions \inserted{(i),} (ii$'$) de la prop 2 montrent que le foncteur ind-représenté par $K$ \uline{dans $C=\mathrm{Cat\,hot}$} \deleted{\unsure{est}} transforme quasi-isomorphismes

% archive sheet number: 44
